\clearpage\thispagestyle{empty}
\pdfbookmark[0]{Quellen zu den Pallas-Berichten vom April, Mai und Juni 1803}{pallas18030406-sources}
\begin{center}\textbf{Quellen zu den Pallas-Berichten vom April, Mai und Juni 1803}\end{center}
\noindent Diese KI-Übertragungen in heutiges Deutsch folgen zwei Auszügen aus den Berichten von Franz Xaver von Zach in \href{https://gdz.sub.uni-goettingen.de/id/PPN236007467}{Gauß’ \textit{Werken}, Band~VI} (Göttingen, 1874), S.~244. Beide Berichte beschreiben Gauß’ Vergleich berechneter Positionen der Pallas mit Beobachtungen von Heinrich Wilhelm Olbers.
\par\smallskip\noindent Für April und Mai geben die \textit{Werke} einen Hinweis zum Aprilbeitrag (XLIV, \href{https://api.digitale-sammlungen.de/iiif/image/v2/bsb10538599_00387/full/1450,/0/default.jpg}{S.~369--375}) in eckigen Klammern und einen Abschnitt aus dem Maibeitrag (LIV, \href{https://api.digitale-sammlungen.de/iiif/image/v2/bsb10538599_00486/full/1450,/0/default.jpg}{S.~466--469}). Der Aprilbericht gibt die \href{https://api.digitale-sammlungen.de/iiif/image/v2/bsb10538599_00392/full/1450,/0/default.jpg}{Elemente VI auf S.~374} an und steht in Zusammenhang mit dem \href{https://api.digitale-sammlungen.de/iiif/image/v2/bsb10538335_00189/full/1450,/0/default.jpg}{Brief an Bode vom 3. März 1803}. Der Maiauszug enthält den \href{https://api.digitale-sammlungen.de/iiif/image/v2/bsb10538599_00488/full/1450,/0/default.jpg}{Vergleich auf S.~468} und die Angabe zur Fortsetzung der \href{https://api.digitale-sammlungen.de/iiif/image/v2/bsb10538598_00597/full/1450,/0/default.jpg}{Ephemeride vom Dezember 1802}. Die Zeitschrift druckt die vollständige Fortsetzung auf \href{https://api.digitale-sammlungen.de/iiif/image/v2/bsb10538599_00489/full/1450,/0/default.jpg}{S.~469}. Der Juniauszug entspricht dem Vergleich auf \href{https://api.digitale-sammlungen.de/iiif/image/v2/bsb10538599_00584/full/1450,/0/default.jpg}{S.~560} am Ende des Beitrags LXV (\href{https://api.digitale-sammlungen.de/iiif/image/v2/bsb10538599_00580/full/1450,/0/default.jpg}{S.~556--560}). Beide Tafeln der \textit{Werke} geben mittlere Bremer Zeiten an.
\par\smallskip\noindent Eine zeitgenössische englische Fassung des Aprilberichts, \href{https://upload.wikimedia.org/wikipedia/commons/d/d6/The_philosophical_magazine._Volume_16%2C_1803._%28IA_s2id13415870%29.pdf#page=53}{„On the new Planet Pallas“, Baron von Zach zugeschrieben}, erschien in \textit{The Philosophical Magazine}, Band~16, Nr.~61 (1803), S.~49--52. Der Beitrag nennt keinen Übersetzer. Grundlage der vorliegenden Übertragungen ist der deutsche Text der \textit{Werke}.
\par\smallskip\noindent Übertragung und Satz wurden mit KI erstellt. Die Digitalisate der deutschen Zeitschrift stammen von der Bayerischen Staatsbibliothek. Das Digitalisat der englischen Zeitschrift stammt aus dem Internet Archive, vermittelt über Wikimedia Commons.
\clearpage
